% Source: https://artofproblemsolving.com/community/c3h2535194_y_in_terms_of_x
% Stable identity: aops-topic-2535194-q01
\textbf{(a)} By definition of the logarithm, $y=\boxed{\log_2x}$; the function $y\mapsto2^y$ is increasing, so this is the only solution.

\textbf{(b)} It is enough to find one $x$ where the formula fails. For $x=2$ the true value is $y=1$, but the formula gives $\sqrt{\dfrac32}+1>2$. (For $x=1$ it gives $2$ instead of $0$, and for $x=4$ it gives $\sqrt{\dfrac52}+1\approx2.58$ instead of $2$.) Indeed $2^y=x$ has the single solution $\log_2x$, and a formula built from square roots cannot agree with it for all $x$: for large $x$, $\sqrt{\dfrac{x+1}{2}}$ grows much faster than $\log_2x$ (for $x=2^{20}$ it exceeds $700$, while $\log_2x=20$).
